\section{Kết luận và Hướng Phát triển}
\label{sec:conclusion}

\subsection{Tổng kết Kết quả Đạt được}
Đồ án đã nghiên cứu và hiện thực thành công một hệ thống mô phỏng và
giải quyết bài toán Bloxorz hoàn chỉnh, đáp ứng toàn diện cả về mặt
kỹ nghệ phần mềm thực hành lẫn phương pháp luận Trí tuệ Nhân tạo thực nghiệm:

\begin{enumerate}[1.]
    \item \textbf{Mô hình hóa Toàn diện Luật chơi Bloxorz:}
    Hệ sinh thái game được xây dựng với đầy đủ các quy tắc vật lý phức tạp:
    động học lăn lật 3 tư thế của khối $1 \times 1 \times 2$, tương tác sàn
    mong manh (gạch cam chịu lực phụ thuộc tư thế), mạng lưới cầu đóng/mở động
    liên kết với nhiều loại công tắc (công tắc tròn, công tắc chữ X), và đặc
    biệt là cơ chế phân thân -- hợp nhất hai khối lập phương độc lập.
    Giao diện 3D trực quan trên nền Ursina Engine kết hợp Tactical Minimap
    mang lại trải nghiệm mượt mà, hỗ trợ quan sát trực tiếp các bước đi do AI
    sinh ra.

    \item \textbf{Hiện thực Chuẩn mực Bộ 4 Giải thuật Tìm kiếm:}
    Xây dựng thành công 4 thuật toán tìm kiếm kinh điển (BFS, DFS, UCS, A*)
    trên nền tảng cấu trúc nút bất biến \texttt{SearchNode}. Kiến trúc chuỗi
    liên kết ngược giúp cắt giảm chi phí cấp phát bộ nhớ cho mỗi nút con từ
    $O(d)$ xuống $O(1)$, loại trừ triệt để lỗi ghi đè tham chiếu trong hàng đợi
    ưu tiên và cho phép tái tạo lộ trình trong thời gian tuyến tính $O(d)$.

    \item \textbf{Thiết kế và Chứng minh Hàm Heuristic:}
    Xây dựng và chứng minh toán học chặt chẽ tính chấp nhận được của hàm
    Heuristic dựa trên khoảng cách Manhattan chia đôi kết hợp phạt định hướng
    sai lệch và chi phí tiếp cận hợp nhất khối phân thân. Chứng minh lý thuyết
    khẳng định rằng A* luôn tìm ra lời giải có tổng chi phí tối ưu toàn cục.

    \item \textbf{Đánh giá Thực nghiệm Định lượng Nghiêm ngặt:}
    Thông qua 40 lượt benchmark lấy mẫu ngẫu nhiên có kiểm soát
    ($\text{seed}=42$), nhóm đã kiểm chứng thực nghiệm sự vượt trội của A*:
    A* cắt tỉa trung bình hơn $56\%$ số nút mở rộng và tăng tốc độ tìm kiếm
    lên tới gần $10$ lần tại các bản đồ phân thân phức tạp so với BFS và UCS.
    Đồng thời, thực nghiệm xác nhận BFS, UCS và A* đạt tính tối ưu tuyệt đối
    $100\%$ về số bước đi, trong khi DFS cho lời giải dài hơn trung bình gấp
    $2,03$ lần.
\end{enumerate}

\subsection{Hướng Phát triển trong Tương lai}
Từ nền tảng của đồ án này, một số hướng nghiên cứu mở rộng đầy tiềm năng
bao gồm:
\begin{itemize}[--]
    \item \textbf{Cơ sở Dữ liệu Mẫu (PDB):}
    Xây dựng các bảng PDB tính toán trước chi phí mở các cụm cầu phức tạp,
    tạo ra các hàm Heuristic có tính tương thích cao hơn và làm tăng đáng kể
    năng lực cắt tỉa tại các mê cung chứa hàng chục công tắc lồng nhau.

    \item \textbf{Thuật toán Tìm kiếm Giới hạn Bộ nhớ:}
    Ứng dụng Iterative Deepening A* (IDA*) hoặc Simplified Memory-Bounded A*
    (SMA*) để giải quyết các bản đồ có quy mô lưới khổng lồ vượt quá giới hạn
    bộ nhớ RAM vật lý.

    \item \textbf{Học Tăng cường Sâu:}
    Huấn luyện mạng nơ-ron tích chập (CNN) kết hợp học tăng cường (DQN / PPO)
    để tác tử tự học các chiến lược định hướng không gian và tương tác công
    tắc từ ảnh chụp ma trận bàn cờ mà không cần phụ thuộc vào hàm Heuristic thủ
    công.
\end{itemize}
